\section{Nov 3} 

	Last time, we talked about the proof of Dudley's theorem with finite sums and then bounding it with an integral. If we have metric space $(T, d)$, with $\mathbb{E} \sup_{t \in T} X_t$. Assume $t \to X_t$ is ``sufficiently continuous.'' 

	\begin{definition}[Guassian Process]
		$(X_t)_{t \in T}$ is a \textbf{Gaussian process (GP)} if 
		\begin{enumerate}
		  \item for every finite $T^\prime \subset T$, $(X_t)_{t \in T^\prime}$ is Gaussian. 
			\item for all $(a_t)_{t \in T^\prime}$, $\sum_{t \in T^\prime} a_t X_t$ is Gaussian. 
		\end{enumerate}
		It is centered if $\mathbb{E}[X_t] = 0$. 
	\end{definition}

	Remark. It is not enough for each $X_t$ is not Gaussian. That is, it is possible to construct where $X_t$ are Gaussians, but their joint distribution is not Gaussian. 

	\begin{lemma}
		A GP is \textbf{fully specified} by a mean function $t \mapsto \mathbb{E}[X_t]$ and a covariance function $\Sigma(s, t) = \Cov(X_t, X_s)$ ($ = \mathbb{E}[X_t X_s]$ for centered GPs). 
	\end{lemma}

	\begin{example}[Finite Gaussian Processes]
	  If the index set is finite, then the whole GP is a giant Gaussian vector 
		\begin{equation}
			(X_t) \sim N(\mu, \Sigma^\ast)
		\end{equation} 
	\end{example}

	\begin{example}
		If $T \subset \mathbb{R}^n$, the \textit{canonical} GP associated with $T$ is 
		\begin{equation}
		  X_t = \langle t, g \rangle, \quad g \sim N(0, I_n)
		\end{equation}
		We can see that the covariance is a dot product since 
		\begin{equation}
			\Sigma(s, t) = \mathbb{E}[X_s, X_t] = \mathbb{E}[(s^T g) (g^T t)] = s^T \mathbb{E}[g g^T] t = s^T t
		\end{equation}
	\end{example}

	GP's are naturally subgaussian w.r.t. a canonical metric on $T$. 
	\begin{align}
		\delta_{X_t - X_s} (\theta) & = \frac{1}{2} \theta^2 \Var(X_s - X_t)  \\ 
																& = \frac{1}{2} \theta^2 \mathbb{E}[|X_s - X_t|^2]
	\end{align}
	So $(X_t)$ is (setting $l=1$) $1$-subgaussian w.r.t. $d(t, s) = \sqrt{\mathbb{E}[|X_s - X_t|^2]} = \sqrt{\Sigma_{ss} + \Sigma_{tt} - 2 \Sigma_{st}}$. 

	\begin{example}
	  If $(X_t)$ are iid $N(0, \sigma^2)$ with $|T| < +\infty$, then 
		\begin{align}
			d(s, t) & = \sqrt{\mathbb{E}[|X_s - X_t|^2]} \\ 
							& = \sqrt{\Var(X_s - X_t)} \\ 
							& = \sqrt{\Var(X_s) + \Var(X_s)} \\ 
							& = \sqrt{2 \sigma^2}
		\end{align}
		for $s \neq t$, and $0$ if $s = t$. Therefore, this is basically a discrete metric. 
	\end{example}

	\begin{example}
		If $T \subset \mathbb{R}^n$ with the canonical GP, then 
		\begin{align}
			d(s, t) & = \sqrt{\Sigma_{ss} + \Sigma_{tt} - 2 \Sigma_{st}} \\ 
							& = \|t - s\|_2
		\end{align}
	\end{example}

	Now we will prove an almost-matching lower bound for GPs. 

	\begin{definition}[Sudakov]
		For a centered GP $(X_t)_{t \in T}$ with canonical metric $d$ on $T$, 
		\begin{equation}
			\mathbb{E}[\sup_{t \in T} X_t] \geq C \sup_{\epsilon > 0} \epsilon \sqrt{\log{N(T, d, \epsilon)}}
		\end{equation}
		for universal constant $C > 0$.\footnote{Note that this doesn't even depend on what the Gaussian process is. It holds for all GPs. }
	\end{definition}

	An equivalent formulation is to look at the discrete sum formulation. 
	\begin{equation}
		 C \sup_{\epsilon > 0} \epsilon \sqrt{\log{N(T, d, \epsilon)}} \leq \mathbb{E}[\sup_{t \in T} X_t] \leq 6 \sum_{k \in \mathbb{Z}} 2^{-k} \sqrt{\log{N(T, d, 2^{-k})}}
	\end{equation}

	So we can use Sudakov in 2 ways. 
	\begin{enumerate}
		\item Lower bound on $N(T, d, \epsilon) \implies$ lower bound on $\mathbb{E}[\sup_t X_t]$. 
		\item Upper bound $\mathbb{E}[\sup_{t \in T} X_t] \implies$ upper bound on $N(T, d, \epsilon)$. 
	\end{enumerate}

	\begin{example}[1st Use Case]
		If $T \subset \mathbb{R}^n$ and $(X_t = \langle t, g \rangle)$ is the canonical GP. Then, 
		\begin{equation}
			N(T, \|\cdot\|_2, \epsilon) \geq \frac{\mathrm{vol}(T)}{\mathrm{vol}(B^n_1)} \epsilon^{-n}
		\end{equation}
		which implies 
		\begin{equation}
			\mathbb{E}[\sup_{t \in T} X_t] \geq C \sup_{\epsilon > 0} \epsilon \sqrt{n \log(1/\epsilon) + \log{\mathrm{vol}(T)} - \log{\mathrm{vol}(B_1^n)}} \sim -\frac{1}{2} n \log{n}
		\end{equation}
		So if $T = [-1, 1]^n$, then $\mathrm{vol}(T) = 2^n$ 
		\begin{equation}
			\implies \mathbb{E}[\|g\|_1] \geq cn
		\end{equation}
	\end{example}
	
	\begin{example}[2nd Use Case]
		Suppose $T$ is the convex hull of finitely many points $\{a_1, \ldots, a_m\}$, i.e. a polytype. An exercise is to show that if $T$ is a convex hull, $\sup_{t \in T} \langle g, t \rangle = \sup_i \langle g, a_i \rangle$, i.e. if you maximize over a convex hull, then it is enough to maximize over the points. Then, 
		\begin{align}
			\mathbb{E}[\sup_{t \in T} \langle t, g \rangle] & = \mathbb{E}[\sup_{i \in [m]}] \underbrace{\langle a_i, g \rangle}_{\|a_i\|_2^2\text{-subgaussian}} \\ 
																											& \leq \sup \|a_i\|_2 \sqrt{2 \log{m}}
		\end{align}
		Using Sudakov, we can see 
		\begin{align}
			N(T, \|\cdot\|_2, \epsilon) \leq m^{C^\prime/\epsilon^2}  
		\end{align}
		So $C^\prime$ is a universal constantthat doesn't depend on the ambient dimension. 
	\end{example}

	\begin{example}
		For canonical GPs for $T \subset \mathbb{R}^n$, we can prove that Dudley is sharp up to $\log{n}$. 
	\end{example}

	\begin{theorem}
		For $T \subset \mathbb{R}^n$, $X_t = \langle t, g \rangle$, 
		\begin{equation}
			c \sup_{\epsilon > 0} \epsilon \sqrt{\log{N(T, \|\cdot\|, \epsilon)}} \leq \mathbb{E}[\sup_t X_t] \leq C (\log{n}) \sup_{\epsilon > 0} \sqrt{\log{N(T, \|\cdot\|_2, \epsilon)}}
		\end{equation} 
	\end{theorem}
	\begin{proof}
	  Exercise.
	\end{proof}

	To prove Sodakov, we need to compare $(X_t)_{t \in T}$ to finitely mand iid Gaussians. More specifically, if we take an $\epsilon$-net $N_{\epsilon} \subset T$, then we need to say that 
	\begin{equation}
		\sup_{t \in N_\epsilon} X_t \approx \sup_{t \in N_\epsilon} Y_t, \quad Y_t \sim N(0, I) \text{ iid}
	\end{equation}


